\subsection{Conjugaison des amplitudes et invariant orienté}
\label{sec:ii09-conjugacion-amplitudes}

La réalisation complexe conserve les trois angles de mélange et
la phase d'interférence comme coordonnées explicites. Écrivons
\(V(\boldsymbol\vartheta,\delta)
=R_{23}U_{13}(\delta)R_{12}\), avec les rotations et la phase
définies dans le développement unitaire. Les arguments des
fonctions trigonométriques sont leurs représentants en radians.

\begin{proposition}[Conjugaison dans la réalisation du mélange]
\label{prop:ii09-conjugacion-unitaria}
Pour les trois angles réels et la phase réelle déclarés,
\[
 V(\boldsymbol\vartheta,-\delta)
       =\overline{V(\boldsymbol\vartheta,\delta)},
 \qquad
 J(\boldsymbol\vartheta,-\delta)
       =-J(\boldsymbol\vartheta,\delta).
\]
L'invariant \(J\) conserve en outre sa valeur sous les changements
diagonaux de phase des bases chargées.
\end{proposition}
\begin{proof}
Les matrices \(R_{12}\) et \(R_{23}\) ont des entrées réelles.
Dans \(U_{13}\), le changement \(\delta\mapsto-\delta\) conjugue
les facteurs exponentiels. La conjugaison du produit donne la
première identité. Dans le quartet
\(V_{11}V_{22}\overline{V_{12}}\overline{V_{21}}\), les
facteurs diagonaux de chaque ligne et colonne s'annulent. La
conjugaison inverse le signe de sa partie imaginaire. On obtient
le résultat de manière équivalente par substitution dans
\[
 J=c_{12}c_{23}c_{13}^{2}s_{12}s_{23}s_{13}\sin\delta.
\]
\end{proof}

Dans le secteur à trois générations, la valeur interne non nulle de
\(J\) exprime l'obstruction à réaliser toutes les amplitudes
par une matrice réelle au moyen de changements diagonaux de phase. C'est
la lecture de violation de charge--parité de la matrice
construite. La lecture des routes conserve, à son tour, l'opération
discrète de conjugaison et de parité ainsi que sa représentation. Leur
compatibilité complète s'exprime par le carré
\[
 \mathcal V(CP_{\rm HMT}\rho)
  =D_u\,\overline{\mathcal V(\rho)}\,D_d^{\dagger},
\]
avec les espaces de représentation et les changements de phase
déclarés. L'identité de conjugaison prouvée ici fixe
exactement l'action que doit transporter ce carré.
